%% kapa: exam=1 round=36 year=2025 period=1 form=A subject=civil_law subject_ko=민법 seq=23 number=23
%% source: Q-Net 감정평가사 기출문제 (cst00309) artlSeq=5231525 file=raw/1008/5231525/2025년 제36회 감정평가사 1차시험 1교시 문제지 원본.pdf
%% conversion: kordoc 4.24.51 -> tools/gen1.py (auto); review=proofread-sonnet; flags=none
\begin{question}{23}
甲은 乙과 乙 소유 X토지에 대한 매매계약을 체결하였다. 그 후 매매대금을 모두 지급한 甲은 X토지를 인도받아 점유ㆍ사용하고 있지만, 乙은 이행기 이후에도 등기를 이전하여 주지 않고 있다. 이에 관한 설명으로 옳지 않은 것은? (다툼이 있으면 판례에 따름)
\begin{choices}
  \item 甲은 乙에 대하여 채권적 등기청구권을 갖는다.
  \item 甲의 乙에 대한 등기청구권의 소멸시효는 진행되지 않는다.
  \item 甲의 乙에 대한 등기청구권은 특별한 사정이 없는 한 그 성질상 양도가 제한된다.
  \item 甲이 X토지를 丙에게 매도하고 인도하였다면, 그 때부터 甲의 乙에 대한 등기청구권의 소멸시효가 진행된다.
  \item 만약 甲이 매매대금을 모두 지급하지 않았다면 甲의 乙에 대한 등기청구권의 소멸시효는 진행되지 않는다.
\end{choices}
\end{question}
\answer{4}
